% ----------------------------------------------------------------------
%
%                            TFMTesis.tex
%
%----------------------------------------------------------------------
%
% Este fichero contiene el "documento maestro" del documento. Lo único
% que hace es configurar el entorno LaTeX e incluir los ficheros .tex
% que contienen cada sección.
%
%----------------------------------------------------------------------
%
% Los ficheros necesarios para este documento son:
%
%       TeXiS/* : ficheros de la plantilla TeXiS.
%       Cascaras/* : ficheros con las partes del documento que no
%          son capítulos ni apéndices (portada, agradecimientos, etc.)
%       Capitulos/*.tex : capítulos de la tesis
%       Apendices/*.tex: apéndices de la tesis
%       constantes.tex: constantes LaTeX
%       config.tex : configuración de la "compilación" del documento
%       guionado.tex : palabras con guiones
%
% Para la bibliografía, además, se necesitan:
%
%       *.bib : ficheros con la información de las referencias
%
% ---------------------------------------------------------------------

\documentclass[12pt,a4paper,twoside]{book}

%
% Definimos  el   comando  \compilaCapitulo,  que   luego  se  utiliza
% (opcionalmente) en config.tex. Quedaría  mejor si también se definiera
% en  ese fichero,  pero por  el modo  en el  que funciona  eso  no es
% posible. Puedes consultar la documentación de ese fichero para tener
% más  información. Definimos también  \compilaApendice, que  tiene el
% mismo  cometido, pero  que se  utiliza para  compilar  únicamente un
% apéndice.
%
%
% Si  queremos   compilar  solo   una  parte  del   documento  podemos
% especificar mediante  \includeonly{...} qué ficheros  son los únicos
% que queremos  que se incluyan.  Esto  es útil por  ejemplo para sólo
% compilar un capítulo.
%
% El problema es que todos aquellos  ficheros que NO estén en la lista
% NO   se  incluirán...  y   eso  también   afecta  a   ficheros  de
% la plantilla...
%
% Total,  que definimos  una constante  con los  ficheros  que siempre
% vamos a querer compilar  (aquellos relacionados con configuración) y
% luego definimos \compilaCapitulo.
\newcommand{\ficherosBasicosTeXiS}{%
TeXiS/TeXiS_pream,TeXiS/TeXiS_cab,TeXiS/TeXiS_bib,TeXiS/TeXiS_cover%
}
\newcommand{\ficherosBasicosTexto}{%
constantes,guionado,Cascaras/bibliografia,config%
}
\newcommand{\compilaCapitulo}[1]{%
\includeonly{\ficherosBasicosTeXiS,\ficherosBasicosTexto,Capitulos/#1}%
}

\newcommand{\compilaApendice}[1]{%
\includeonly{\ficherosBasicosTeXiS,\ficherosBasicosTexto,Apendices/#1}%
}

%- - - - - - - - - - - - - - - - - - - - - - - - - - - - - - - - - - -
%            Preámbulo del documento. Configuraciones varias
%- - - - - - - - - - - - - - - - - - - - - - - - - - - - - - - - - - -

% Define  el  tipo  de  compilación que  estamos  haciendo.   Contiene
% definiciones  de  constantes que  cambian  el  comportamiento de  la
% compilación. Debe incluirse antes del paquete TeXiS/TeXiS.sty
%---------------------------------------------------------------------
%
%                          config.tex
%
%---------------------------------------------------------------------
%
% Contiene la  definición de constantes  que determinan el modo  en el
% que se compilará el documento.
%
%---------------------------------------------------------------------
%
% En concreto, podemos  indicar si queremos "modo release",  en el que
% no  aparecerán  los  comentarios  (creados  mediante  \com{Texto}  o
% \comp{Texto}) ni los "por  hacer" (creados mediante \todo{Texto}), y
% sí aparecerán los índices. El modo "debug" (o mejor dicho en modo no
% "release" muestra los índices  (construirlos lleva tiempo y son poco
% útiles  salvo  para   la  versión  final),  pero  sí   el  resto  de
% anotaciones.
%
% Si se compila con LaTeX (no  con pdflatex) en modo Debug, también se
% muestran en una esquina de cada página las entradas (en el índice de
% palabras) que referencian  a dicha página (consulta TeXiS_pream.tex,
% en la parte referente a show).
%
% El soporte para  el índice de palabras en  TeXiS es embrionario, por
% lo  que no  asumas que  esto funcionará  correctamente.  Consulta la
% documentación al respecto en TeXiS_pream.tex.
%
%
% También  aquí configuramos  si queremos  o  no que  se incluyan  los
% acrónimos  en el  documento final  en la  versión release.  Para eso
% define (o no) la constante \acronimosEnRelease.
%
% Utilizando \compilaCapitulo{nombre}  podemos también especificar qué
% capítulo(s) queremos que se compilen. Si no se pone nada, se compila
% el documento  completo.  Si se pone, por  ejemplo, 01Introduccion se
% compilará únicamente el fichero Capitulos/01Introduccion.tex
%
% Para compilar varios  capítulos, se separan sus nombres  con comas y
% no se ponen espacios de separación.
%
% En realidad  la macro \compilaCapitulo  está definida en  el fichero
% principal tesis.tex.
%
%---------------------------------------------------------------------


% Comentar la línea si no se compila en modo release.
% TeXiS hará el resto.
% ¡¡¡Si cambias esto, haz un make clean antes de recompilar!!!
\def\release{1}


% Descomentar la linea si se quieren incluir los
% acrónimos en modo release (en modo debug
% no se incluirán nunca).
% ¡¡¡Si cambias esto, haz un make clean antes de recompilar!!!
%\def\acronimosEnRelease{1}


% Descomentar la línea para establecer el capítulo que queremos
% compilar

% \compilaCapitulo{01Introduccion}
% \compilaCapitulo{02EstructuraYGeneracion}
% \compilaCapitulo{03Edicion}
% \compilaCapitulo{04Imagenes}
% \compilaCapitulo{05Bibliografia}
% \compilaCapitulo{06Makefile}

% \compilaApendice{01AsiSeHizo}

% Variable local para emacs, para  que encuentre el fichero maestro de
% compilación y funcionen mejor algunas teclas rápidas de AucTeX
%%%
%%% Local Variables:
%%% mode: latex
%%% TeX-master: "./Tesis.tex"
%%% End:


% Paquete de la plantilla
\usepackage{TeXiS/TeXiS}

% Incluimos el fichero con comandos de constantes
%---------------------------------------------------------------------
%
%                          constantes.tex
%
%---------------------------------------------------------------------
%
% Fichero que  declara nuevos comandos LaTeX  sencillos realizados por
% comodidad en la escritura de determinadas palabras
%
%---------------------------------------------------------------------

%%%%%%%%%%%%%%%%%%%%%%%%%%%%%%%%%%%%%%%%%%%%%%%%%%%%%%%%%%%%%%%%%%%%%%
% Comando: 
%
%       \titulo
%
% Resultado: 
%
% Escribe el título del documento.
%%%%%%%%%%%%%%%%%%%%%%%%%%%%%%%%%%%%%%%%%%%%%%%%%%%%%%%%%%%%%%%%%%%%%%
\def\titulo{\textsc{TeXiS}: Una plantilla de \LaTeX\
  para Tesis y otros documentos}

%%%%%%%%%%%%%%%%%%%%%%%%%%%%%%%%%%%%%%%%%%%%%%%%%%%%%%%%%%%%%%%%%%%%%%
% Comando: 
%
%       \autor
%
% Resultado: 
%
% Escribe el autor del documento.
%%%%%%%%%%%%%%%%%%%%%%%%%%%%%%%%%%%%%%%%%%%%%%%%%%%%%%%%%%%%%%%%%%%%%%
\def\autor{Marco Antonio y Pedro Pablo G\'omez Mart\'in}

% Variable local para emacs, para  que encuentre el fichero maestro de
% compilación y funcionen mejor algunas teclas rápidas de AucTeX

%%%
%%% Local Variables:
%%% mode: latex
%%% TeX-master: "tesis.tex"
%%% End:


% Sacamos en el log de la compilación el copyright
%\typeout{Copyright Marco Antonio and Pedro Pablo Gomez Martin}

%
% "Metadatos" para el PDF
%
\ifpdf\hypersetup{%
    pdftitle = {\titulo},
    pdfsubject = {Trabajo de Fin de Grado},
    pdfkeywords = {Trabajo de Fin de Grado, RPG, motor de videojuegos, videojuegos, editor de videojuegos},
    pdfauthor = {\textcopyright\ \autor},
    pdfcreator = {\LaTeX\ con el paquete \flqq hyperref\frqq},
    pdfproducer = {pdfeTeX-0.\the\pdftexversion\pdftexrevision},
    }
    \pdfinfo{/CreationDate (\today)}
\fi


%- - - - - - - - - - - - - - - - - - - - - - - - - - - - - - - - - - -
%                        Documento
%- - - - - - - - - - - - - - - - - - - - - - - - - - - - - - - - - - -
\begin{document}

% Incluimos el  fichero de definición de guionado  de algunas palabras
% que LaTeX no ha dividido como debería
%----------------------------------------------------------------
%
%                          guionado.tex
%
%----------------------------------------------------------------
%
% Fichero con algunas divisiones de palabras que LaTeX no
% hace correctamente si no se le da alguna ayuda.
%
%----------------------------------------------------------------

\hyphenation{
% a
abs-trac-to
abs-trac-tos
abs-trac-ta
abs-trac-tas
ac-tua-do-res
a-gra-de-ci-mien-tos
ana-li-za-dor
an-te-rio-res
an-te-rior-men-te
apa-rien-cia
a-pro-pia-do
a-pro-pia-dos
a-pro-pia-da
a-pro-pia-das
a-pro-ve-cha-mien-to
a-que-llo
a-que-llos
a-que-lla
a-que-llas
a-sig-na-tu-ra
a-sig-na-tu-ras
a-so-cia-da
a-so-cia-das
a-so-cia-do
a-so-cia-dos
au-to-ma-ti-za-do
% b
batch
bi-blio-gra-fía
bi-blio-grá-fi-cas
bien
bo-rra-dor
boo-l-ean-expr
% c
ca-be-ce-ra
call-me-thod-ins-truc-tion
cas-te-lla-no
cir-cuns-tan-cia
cir-cuns-tan-cias
co-he-ren-te
co-he-ren-tes
co-he-ren-cia
co-li-bri
co-men-ta-rio
co-mer-cia-les
co-no-ci-mien-to
cons-cien-te
con-si-de-ra-ba
con-si-de-ra-mos
con-si-de-rar-se
cons-tan-te
cons-trucción
cons-tru-ye
cons-tru-ir-se
con-tro-le
co-rrec-ta-men-te
co-rres-pon-den
co-rres-pon-dien-te
co-rres-pon-dien-tes
co-ti-dia-na
co-ti-dia-no
crean
cris-ta-li-zan
cu-rri-cu-la
cu-rri-cu-lum
cu-rri-cu-lar
cu-rri-cu-la-res
% d
de-di-ca-do
de-di-ca-dos
de-di-ca-da
de-di-ca-das
de-rro-te-ro
de-rro-te-ros
de-sa-rro-llo
de-sa-rro-llos
de-sa-rro-lla-do
de-sa-rro-lla-dos
de-sa-rro-lla-da
de-sa-rro-lla-das
de-sa-rro-lla-dor
de-sa-rro-llar
des-cri-bi-re-mos
des-crip-ción
des-crip-cio-nes
des-cri-to
des-pués
de-ta-lla-do
de-ta-lla-dos
de-ta-lla-da
de-ta-lla-das
di-a-gra-ma
di-a-gra-mas
di-se-ños
dis-po-ner
dis-po-ni-bi-li-dad
do-cu-men-ta-da
do-cu-men-to
do-cu-men-tos
% e
edi-ta-do
e-du-ca-ti-vo
e-du-ca-ti-vos
e-du-ca-ti-va
e-du-ca-ti-vas
e-la-bo-ra-do
e-la-bo-ra-dos
e-la-bo-ra-da
e-la-bo-ra-das
es-co-llo
es-co-llos
es-tu-dia-do
es-tu-dia-dos
es-tu-dia-da
es-tu-dia-das
es-tu-dian-te
e-va-lua-cio-nes
e-va-lua-do-res
exis-ten-tes
exhaus-ti-va
ex-pe-rien-cia
ex-pe-rien-cias
% f
for-ma-li-za-do
% g
ge-ne-ra-ción
ge-ne-ra-dor
ge-ne-ra-do-res
ge-ne-ran
% h
he-rra-mien-ta
he-rra-mien-tas
% i
i-dio-ma
i-dio-mas
im-pres-cin-di-ble
im-pres-cin-di-bles
in-de-xa-do
in-de-xa-dos
in-de-xa-da
in-de-xa-das
in-di-vi-dual
in-fe-ren-cia
in-fe-ren-cias
in-for-ma-ti-ca
in-gre-dien-te
in-gre-dien-tes
in-me-dia-ta-men-te
ins-ta-la-do
ins-tan-cias
in-te-rac-ción
% j
% k
% l
len-gua-je
li-be-ra-to-rio
li-be-ra-to-rios
li-be-ra-to-ria
li-be-ra-to-rias
li-mi-ta-do
li-te-ra-rio
li-te-ra-rios
li-te-ra-ria
li-te-ra-rias
lo-tes
% m
ma-ne-ra
ma-nual
mas-que-ra-de
ma-yor
me-mo-ria
mi-nis-te-rio
mi-nis-te-rios
mo-de-lo
mo-de-los
mo-de-la-do
mo-du-la-ri-dad
mo-vi-mien-to
% n
na-tu-ral
ni-vel
nues-tro
% o
obs-tan-te
o-rien-ta-do
o-rien-ta-dos
o-rien-ta-da
o-rien-ta-das
% p
pa-ra-le-lo
pa-ra-le-la
par-ti-cu-lar
par-ti-cu-lar-men-te
pe-da-gó-gi-ca
pe-da-gó-gi-cas
pe-da-gó-gi-co
pe-da-gó-gi-cos
pe-rio-di-ci-dad
per-so-na-je
plan-te-a-mien-to
plan-te-a-mien-tos
po-si-ción
pre-fe-ren-cia
pre-fe-ren-cias
pres-cin-di-ble
pres-cin-di-bles
pri-me-ra
pro-ble-ma
pro-ble-mas
pró-xi-mo
pu-bli-ca-cio-nes
pu-bli-ca-do
% q
% r
rá-pi-da
rá-pi-do
ra-zo-na-mien-to
ra-zo-na-mien-tos
re-a-li-zan-do
re-fe-ren-cia
re-fe-ren-cias
re-fe-ren-cia-da
re-fe-ren-cian
re-le-van-tes
re-pre-sen-ta-do
re-pre-sen-ta-dos
re-pre-sen-ta-da
re-pre-sen-ta-das
re-pre-sen-tar-lo
re-qui-si-to
re-qui-si-tos
res-pon-der
res-pon-sa-ble
% s
se-pa-ra-do
si-guien-do
si-guien-te
si-guien-tes
si-guie-ron
si-mi-lar
si-mi-la-res
si-tua-ción
% t
tem-pe-ra-ments
te-ner
trans-fe-ren-cia
trans-fe-ren-cias
% u
u-sua-rio
Unreal-Ed
% v
va-lor
va-lo-res
va-rian-te
ver-da-de-ro
ver-da-de-ros
ver-da-de-ra
ver-da-de-ras
ver-da-de-ra-men-te
ve-ri-fi-ca
% w
% x
% y
% z
}
% Variable local para emacs, para que encuentre el fichero
% maestro de compilación
%%%
%%% Local Variables:
%%% mode: latex
%%% TeX-master: "./Tesis.tex"
%%% End:


% Marcamos  el inicio  del  documento para  la  numeración de  páginas
% (usando números romanos para esta primera fase).
\frontmatter
\pagestyle{empty}

%---------------------------------------------------------------------
%
%                          configCover.tex
%
%---------------------------------------------------------------------
%
% cover.tex
% Copyright 2009 Marco Antonio Gomez-Martin, Pedro Pablo Gomez-Martin
%
% This file belongs to the TeXiS manual, a LaTeX template for writting
% Thesis and other documents. The complete last TeXiS package can
% be obtained from http://gaia.fdi.ucm.es/projects/texis/
%
% Although the TeXiS template itself is distributed under the 
% conditions of the LaTeX Project Public License
% (http://www.latex-project.org/lppl.txt), the manual content
% uses the CC-BY-SA license that stays that you are free:
%
%    - to share & to copy, distribute and transmit the work
%    - to remix and to adapt the work
%
% under the following conditions:
%
%    - Attribution: you must attribute the work in the manner
%      specified by the author or licensor (but not in any way that
%      suggests that they endorse you or your use of the work).
%    - Share Alike: if you alter, transform, or build upon this
%      work, you may distribute the resulting work only under the
%      same, similar or a compatible license.
%
% The complete license is available in
% http://creativecommons.org/licenses/by-sa/3.0/legalcode
%
%---------------------------------------------------------------------
%
% Fichero que contiene la configuración de la portada y de la 
% primera hoja del documento.
%
%---------------------------------------------------------------------


% Pueden configurarse todos los elementos del contenido de la portada
% utilizando comandos.

%%%%%%%%%%%%%%%%%%%%%%%%%%%%%%%%%%%%%%%%%%%%%%%%%%%%%%%%%%%%%%%%%%%%%%
% Título del documento:
% \tituloPortada{titulo}
% Nota:
% Si no se define se utiliza el del \titulo. Este comando permite
% cambiar el título de forma que se especifiquen dónde se quieren
% los retornos de carro cuando se utilizan fuentes grandes.
%%%%%%%%%%%%%%%%%%%%%%%%%%%%%%%%%%%%%%%%%%%%%%%%%%%%%%%%%%%%%%%%%%%%%%
\tituloPortada{%
Motor y editor de videojuegos en 2D enfocado al desarrollo de juegos RPG
}


%%%%%%%%%%%%%%%%%%%%%%%%%%%%%%%%%%%%%%%%%%%%%%%%%%%%%%%%%%%%%%%%%%%%%%
% Título del documento en inglés:
% \tituloPortadaEng{titulo}
% Nota:
% Si no se define se utiliza el del \titulo. Este comando permite
% cambiar el título de forma que se especifiquen dónde se quieren
% los retornos de carro cuando se utilizan fuentes grandes.
%%%%%%%%%%%%%%%%%%%%%%%%%%%%%%%%%%%%%%%%%%%%%%%%%%%%%%%%%%%%%%%%%%%%%%
\tituloPortadaEng{%
2D video game engine and editor focused on RPG game development
}

%%%%%%%%%%%%%%%%%%%%%%%%%%%%%%%%%%%%%%%%%%%%%%%%%%%%%%%%%%%%%%%%%%%%%%
% Autor del documento:
% \autorPortada{Nombre}
% Se utiliza en la portada y en el valor por defecto del
% primer subtítulo de la segunda portada.
%%%%%%%%%%%%%%%%%%%%%%%%%%%%%%%%%%%%%%%%%%%%%%%%%%%%%%%%%%%%%%%%%%%%%%
\autorPortada{Miguel Curros Garc\'ia\\Alejandro Gonz\'alez S\'anchez\\Alejandro Mass\'o Mart\'inez}

%%%%%%%%%%%%%%%%%%%%%%%%%%%%%%%%%%%%%%%%%%%%%%%%%%%%%%%%%%%%%%%%%%%%%%
% Fecha de publicación:
% \fechaPublicacion{Fecha}
% Puede ser vacío. Aparece en la última línea de ambas portadas
%%%%%%%%%%%%%%%%%%%%%%%%%%%%%%%%%%%%%%%%%%%%%%%%%%%%%%%%%%%%%%%%%%%%%%
% Descomentar para que ponga siempre la fecha actual
%\fechaPublicacion{\today}
\fechaPublicacion{26 de mayo de 2025}

%%%%%%%%%%%%%%%%%%%%%%%%%%%%%%%%%%%%%%%%%%%%%%%%%%%%%%%%%%%%%%%%%%%%%%
% Imagen de la portada (y escala)
% \imagenPortada{Fichero}
% \escalaImagenPortada{Numero}
% Si no se especifica, se utiliza la imagen TODO.pdf
%%%%%%%%%%%%%%%%%%%%%%%%%%%%%%%%%%%%%%%%%%%%%%%%%%%%%%%%%%%%%%%%%%%%%%
% imagen en blanco y negro
%\imagenPortada{Imagenes/Vectorial/escudoUCM}
%imagen en color
\imagenPortada{Imagenes/Bitmap/escudoUCMcolor}
\escalaImagenPortada{.2}

%%%%%%%%%%%%%%%%%%%%%%%%%%%%%%%%%%%%%%%%%%%%%%%%%%%%%%%%%%%%%%%%%%%%%%
% Tipo de documento.
% \tipoDocumento{Tipo}
% Para el texto justo debajo del escudo.
% Si no se indica, se utiliza "TESIS DOCTORAL".
%%%%%%%%%%%%%%%%%%%%%%%%%%%%%%%%%%%%%%%%%%%%%%%%%%%%%%%%%%%%%%%%%%%%%%
\tipoDocumento{Trabajo de Fin de Grado}

%%%%%%%%%%%%%%%%%%%%%%%%%%%%%%%%%%%%%%%%%%%%%%%%%%%%%%%%%%%%%%%%%%%%%%
% Institución/departamento asociado al documento.
% \institucion{Nombre}
% Puede tener varias líneas. Se utiliza en las dos portadas.
% Si no se indica aparecerá vacío.
%%%%%%%%%%%%%%%%%%%%%%%%%%%%%%%%%%%%%%%%%%%%%%%%%%%%%%%%%%%%%%%%%%%%%%
\institucion{%
Grado en Desarrollo de Videojuegos\\[0.2em]
Facultad de Inform\'atica\\[0.2em]
Universidad Complutense de Madrid
}

%%%%%%%%%%%%%%%%%%%%%%%%%%%%%%%%%%%%%%%%%%%%%%%%%%%%%%%%%%%%%%%%%%%%%%
% Director del trabajo.
% \directorPortada{Nombre}
% Se utiliza para el valor por defecto del segundo subtítulo, donde
% se indica quién es el director del trabajo.
% Si se fuerza un subtítulo distinto, no hace falta definirlo.
%%%%%%%%%%%%%%%%%%%%%%%%%%%%%%%%%%%%%%%%%%%%%%%%%%%%%%%%%%%%%%%%%%%%%%
\directorPortada{Pedro Pablo G\'omez Mart\'in}


%%%%%%%%%%%%%%%%%%%%%%%%%%%%%%%%%%%%%%%%%%%%%%%%%%%%%%%%%%%%%%%%%%%%%%
% Colaborador en la dirección del trabajo.
% \colaboradorPortada{Nombre}
% Se utiliza para el valor por defecto del segundo subtítulo, donde
% se indica quién es el colaborador en la dirección del trabajo.
% Si se fuerza un subtítulo distinto, no hace falta definirlo.
%%%%%%%%%%%%%%%%%%%%%%%%%%%%%%%%%%%%%%%%%%%%%%%%%%%%%%%%%%%%%%%%%%%%%%
%\colaboradorPortada{\textcolor{red}{Colaborador 1\\Colaborador 2}}


%%%%%%%%%%%%%%%%%%%%%%%%%%%%%%%%%%%%%%%%%%%%%%%%%%%%%%%%%%%%%%%%%%%%%%
% Texto del primer subtítulo de la segunda portada.
% \textoPrimerSubtituloPortada{Texto}
% Para configurar el primer "texto libre" de la segunda portada.
% Si no se especifica se indica "Memoria que presenta para optar al
% título de Doctor en Informática" seguido del \autorPortada.
%%%%%%%%%%%%%%%%%%%%%%%%%%%%%%%%%%%%%%%%%%%%%%%%%%%%%%%%%%%%%%%%%%%%%%
\textoPrimerSubtituloPortada{%
\textbf{Trabajo de Fin de Grado en Desarrollo de Videojuegos}\\ [0.3em]
}

%%%%%%%%%%%%%%%%%%%%%%%%%%%%%%%%%%%%%%%%%%%%%%%%%%%%%%%%%%%%%%%%%%%%%%
% Texto del segundo subtítulo de la segunda portada.
% \textoSegundoSubtituloPortada{Texto}
% Para configurar el segundo "texto libre" de la segunda portada.
% Si no se especifica se indica "Dirigida por el Doctor" seguido
% del \directorPortada.
%%%%%%%%%%%%%%%%%%%%%%%%%%%%%%%%%%%%%%%%%%%%%%%%%%%%%%%%%%%%%%%%%%%%%%
\textoSegundoSubtituloPortada{%
\textbf{Convocatoria: }\textit{Junio \the\year}%\\[0.2em]
%\textbf{Calificación: }\textit{\textcolor{red}{Nota}}
}

%%%%%%%%%%%%%%%%%%%%%%%%%%%%%%%%%%%%%%%%%%%%%%%%%%%%%%%%%%%%%%%%%%%%%%
% \explicacionDobleCara
% Si se utiliza, se aclara que el documento está preparado para la
% impresión a doble cara.
%%%%%%%%%%%%%%%%%%%%%%%%%%%%%%%%%%%%%%%%%%%%%%%%%%%%%%%%%%%%%%%%%%%%%%
\explicacionDobleCara

%%%%%%%%%%%%%%%%%%%%%%%%%%%%%%%%%%%%%%%%%%%%%%%%%%%%%%%%%%%%%%%%%%%%%%
% \isbn
% Si se utiliza, aparecerá el ISBN detrás de la segunda portada.
%%%%%%%%%%%%%%%%%%%%%%%%%%%%%%%%%%%%%%%%%%%%%%%%%%%%%%%%%%%%%%%%%%%%%%
%\isbn{978-84-692-7109-4}


%%%%%%%%%%%%%%%%%%%%%%%%%%%%%%%%%%%%%%%%%%%%%%%%%%%%%%%%%%%%%%%%%%%%%%
% \copyrightInfo
% Si se utiliza, aparecerá información de los derechos de copyright
% detrás de la segunda portada.
%%%%%%%%%%%%%%%%%%%%%%%%%%%%%%%%%%%%%%%%%%%%%%%%%%%%%%%%%%%%%%%%%%%%%%
%\copyrightInfo{\autor}


%%
%% Creamos las portadas
%%
\makeCover

% Variable local para emacs, para que encuentre el fichero
% maestro de compilación
%%%
%%% Local Variables:
%%% mode: latex
%%% TeX-master: "../Tesis.tex"
%%% End:

%\include{Cascaras/autorizacion}
% +--------------------------------------------------------------------+
% | Dedication Page (Optional)
% +--------------------------------------------------------------------+

\chapter*{Dedicatoria}

\begin{flushright}
\begin{minipage}[c]{8.5cm}
\flushright{\textit{Dedicatoria de Miguel. \textbf{- Miguel}}}

\bigskip

\flushright{\textit{Dedicatoria de Alex. \textbf{- Alex G.}}}

\bigskip

\flushright{\textit{Dedicatoria de Alex. \textbf{- Alex M.}}}
\end{minipage}
\end{flushright}
% +--------------------------------------------------------------------+
% | Acknowledgements Page (Optional)                                   |
% +--------------------------------------------------------------------+

\chapter*{Agradecimientos}

Gracias, de nuevo, a todos nuestros familiares, por haber sido la guía \comillas{detrás de las sombras} todo este tiempo. Sin vosotros, nada de esto hubiese sido posible.

\medskip

Gracias, también, a todos nuestros amigos que decidieron dedicar parte de su valioso tiempo para poder probar nuestras herramientas y darnos sus opiniones. Sin vosotros, la entrega final no hubiese sido posible. 

\medskip

Gracias a todos nuestros profesores, tanto del Grado, como del Colegio e Instituto por habernos brindado la oportunidad de estar aquí haciendo lo que más nos gusta, y por habernos enseñado tantas valiosas lecciones no siempre relacionadas con los estudios.

\medskip

Y por último, y no por ello menos importante, gracias a Pedro Pablo, nuestro tutor, por habernos brindado su confianza desde el primer día y haber creído en nosotros para poder acabar esto. Gracias por habernos tenido tanta paciencia, que a veces era muy necesaria, y gracias por habernos dado los tirones de oreja cuando tocaba.

\medskip

Este trabajo va por todos vosotros.












\chapter*{Resumen}

\section*{\tituloPortadaVal}

Un resumen en castellano de media página, incluyendo el título en castellano. A continuación, se escribirá una lista de no más de 10 palabras clave.


\section*{Palabras clave}
   
\noindent Máximo 10 palabras clave separadas por comas

   



\begin{otherlanguage}{english}
\chapter*{Abstract}

\section*{\tituloPortadaEngVal}

Although role-playing video games (RPGs) are among the most demanded genres in today's industry, specific engines and editors designed for these type of games are scarce. The most used ones tend to be too generic and are not focused on RPGs specifically.

\smallskip

In addition to this, the available open source RPG engines and editors are often not very user-friendly for newcomers, as they use mechanisms and structures designed for users with advanced knowledge. Those that do offer a more straightforward and practical interface are often locked behind a paywall, so amateur designers find themselves spending a large amount of money on the license for a software they will rarely use.

\medskip

To address these problems, this project presents the development of a 2D-video game engine and editor specifically oriented towards RPGs. The engine offers cross-platform support for Windows, MacOS, Linux, and Android, while the editor is available for Windows, MacOS, and Linux. Its goal is to ease the creation of this type of games for users without experience in development or programming, while maintaning more advanced features so that experienced users can generate more complete games, all using open-source and freely accessible tools.


\section*{Keywords}

\noindent Role-playing Video Games, Video Game Development, Video Game Editor, Development Tool, Video Game Engine.




% Si el trabajo se escribe en inglés, comentar esta línea y descomentar
% otra igual que hay justo antes de \end{document}
\end{otherlanguage}

\ifx\generatoc\undefined
\else
%---------------------------------------------------------------------
%
%                          TeXiS_toc.tex
%
%---------------------------------------------------------------------
%
% TeXiS_toc.tex
% Copyright 2009 Marco Antonio Gomez-Martin, Pedro Pablo Gomez-Martin
%
% This file belongs to TeXiS, a LaTeX template for writting
% Thesis and other documents. The complete last TeXiS package can
% be obtained from http://gaia.fdi.ucm.es/projects/texis/
%
% This work may be distributed and/or modified under the
% conditions of the LaTeX Project Public License, either version 1.3
% of this license or (at your option) any later version.
% The latest version of this license is in
%   http://www.latex-project.org/lppl.txt
% and version 1.3 or later is part of all distributions of LaTeX
% version 2005/12/01 or later.
%
% This work has the LPPL maintenance status `maintained'.
% 
% The Current Maintainers of this work are Marco Antonio Gomez-Martin
% and Pedro Pablo Gomez-Martin
%
%---------------------------------------------------------------------
%
% Contiene  los  comandos  para  generar los  índices  del  documento,
% entendiendo por índices las tablas de contenidos.
%
% Genera  el  índice normal  ("tabla  de  contenidos"),  el índice  de
% figuras y el de tablas. También  crea "marcadores" en el caso de que
% se esté compilando con pdflatex para que aparezcan en el PDF.
%
%---------------------------------------------------------------------


% Primero un poquito de configuración...


% Pedimos que inserte todos los epígrafes hasta el nivel \subsection en
% la tabla de contenidos.
\setcounter{tocdepth}{2} 

% Le  pedimos  que nos  numere  todos  los  epígrafes hasta  el  nivel
% \subsubsection en el cuerpo del documento.
\setcounter{secnumdepth}{3} 


% Creamos los diferentes índices.

% Lo primero un  poco de trabajo en los marcadores  del PDF. No quiero
% que  salga una  entrada  por cada  índice  a nivel  0...  si no  que
% aparezca un marcador "Índices", que  tenga dentro los otros tipos de
% índices.  Total, que creamos el marcador "Índices".
% Antes de  la creación  de los índices,  se añaden los  marcadores de
% nivel 1.

\ifpdf
   \pdfbookmark{Índices}{indices}
\fi

% Tabla de contenidos.
%
% La  inclusión  de '\tableofcontents'  significa  que  en la  primera
% pasada  de  LaTeX  se  crea   un  fichero  con  extensión  .toc  con
% información sobre la tabla de contenidos (es conceptualmente similar
% al  .bbl de  BibTeX, creo).  En la  segunda ejecución  de  LaTeX ese
% documento se utiliza para  generar la verdadera página de contenidos
% usando la  información sobre los  capítulos y demás guardadas  en el
% .toc
\ifpdf
   \pdfbookmark[1]{Tabla de Contenidos}{tabla de contenidos}
\fi

\cabeceraEspecial{\'Indice}

\tableofcontents

\newpage 

% Índice de figuras
%
% La idea es semejante que para  el .toc del índice, pero ahora se usa
% extensión .lof (List Of Figures) con la información de las figuras.

\ifpdf
   \pdfbookmark[1]{Índice de figuras}{indice de figuras}
\fi

\cabeceraEspecial{\'Indice de figuras}

\listoffigures

\newpage

% Índice de tablas
% Como antes, pero ahora .lot (List Of Tables)

\ifpdf
   \pdfbookmark[1]{Índice de tablas}{indice de tablas}
\fi

\cabeceraEspecial{\'Indice de tablas}

\listoftables

\newpage

% Variable local para emacs, para  que encuentre el fichero maestro de
% compilación y funcionen mejor algunas teclas rápidas de AucTeX

%%%
%%% Local Variables:
%%% mode: latex
%%% TeX-master: "../Tesis.tex"
%%% End:

\fi

% Marcamos el  comienzo de  los capítulos (para  la numeración  de las
% páginas) y ponemos la cabecera normal
\mainmatter

\pagestyle{fancy}
\restauraCabecera


\chapter{Introducción}
\label{cap:introduccion}

\chapterquote{Frase célebre dicha por alguien inteligente}{Autor}

Según la normativa  para Trabajos de Fin de Grado\footnote{\url{https://informatica.ucm.es/file/normativatfg-2021-2022?ver} (ver versión actualizada para cada curso académico)}, la memoria incluirá una portada normalizada con la siguiente información: título en castellano, título en inglés, autores, profesor director, codirector si es el caso, curso
académico e identificación de la asignatura (Trabajo de fin de grado del Grado en -
nombre del grado correspondiente-, Facultad de Informática, Universidad Complutense de Madrid). Los datos referentes al título y director (y codirector en su caso) deben corresponder a los publicados en la página web de TFG.

La memoria debe incluir la descripción detallada de la propuesta hardware/software realizada y ha de contener:
\begin{enumerate}
\renewcommand{\theenumi}{\alph{enumi}}
	\item un índice,
	\item un resumen y una lista de no más de 10 palabras clave para su búsqueda bibliográfica, ambos en castellano e inglés,
	\item una introducción con los antecedentes, objetivos y plan de trabajo,
	\item resultados y discusión crítica y razonada de los mismos, con sus conclusiones,
	\item bibliografía.
\end{enumerate}
Para facilitar la escritura de la memoria siguiendo esta estructura, el estudiante podrá usar las plantillas en LaTeX o Word preparadas al efecto y publicadas en la página web de TFG.

La memoria constará de un mínimo de 25 páginas para los proyectos realizados por un único estudiante, y de al menos 5 páginas más por cada integrante adicional del grupo. En este número de páginas solo se tiene en cuenta el contenido correspondiente a los apartados c y d del punto anterior.

La memoria puede estar escrita en castellano o inglés, pero en el primer caso la introducción y las conclusiones deben aparecer también en inglés. Las memorias de los TFG matriculados en el grupo I deberán estar escritas íntegramente en inglés, excepto por lo especificado en los puntos 1 y 2 anteriores (título, resumen y lista de palabras clave).

En caso de trabajos no unipersonales, cada participante indicará en la memoria su contribución al proyecto con una extensión de al menos dos páginas por cada uno de los participantes.

Todo el material no original, ya sea texto o figuras, deberá ser convenientemente citado y referenciado. En el caso de material complementario se deben respetar las licencias y \emph{copyrights} asociados al software y hardware que se emplee. En caso contrario no se autorizará la defensa, sin menoscabo de otras acciones que correspondan.


\section{Motivación}
Introducción al tema del TFG.


\section{Objetivos}
Descripción de los objetivos del trabajo.


\section{Plan de trabajo}
Aquí se describe el plan de trabajo a seguir para la consecución de los objetivos descritos en el apartado anterior.



\section{Explicaciones adicionales sobre el uso de esta plantilla}
Si quieres cambiar el \textbf{estilo del título} de los capítulos del documento, edita el fichero \verb|TeXiS\TeXiS_pream.tex| y comenta la línea \verb|\usepackage[Lenny]{fncychap}| para dejar el estilo básico de \LaTeX.

Si no te gusta que no haya \textbf{espacios entre párrafos} y quieres dejar un pequeño espacio en blanco, no metas saltos de línea (\verb|\\|) al final de los párrafos. En su lugar, busca el comando  \verb|\setlength{\parskip}{0.2ex}| en \verb|TeXiS\TeXiS_pream.tex| y aumenta el valor de $0.2ex$ a, por ejemplo, $1ex$.

TFGTeXiS se ha elaborado a partir de la plantilla de TeXiS\footnote{\url{http://gaia.fdi.ucm.es/research/texis/}}, creada por Marco Antonio y Pedro Pablo Gómez Martín para escribir su tesis doctoral. Para explicaciones más extensas y detalladas sobre cómo usar esta plantilla, recomendamos la lectura del documento \texttt{TeXiS-Manual-1.0.pdf} que acompaña a esta plantilla.

El siguiente texto se genera con el comando \verb|\lipsum[2-20]| que viene a continuación en el fichero .tex. El único propósito es mostrar el aspecto de las páginas usando esta plantilla. Quita este comando y, si quieres, comenta o elimina el paquete \textit{lipsum} al final de \verb|TeXiS\TeXiS_pream.tex|

\subsection{Texto de prueba}


\lipsum[2-20]
\chapter{Estado de la cuestión}
\label{cap:estadoDeLaCuestion}

\begin{resumen}
En este capítulo vamos a hablar sobre los videojuegos de rol, introducir un poco su historia, hablar sobre el desarrollo de videojuegos, qué es un motor, para qué sirve, qué lo compone, qué es un editor, para qué sirve, y cuáles son los editores especializados en el desarrollo de videojuegos de rol.
\end{resumen}

\section{Sobre los videojuegos de rol}
Para poder definir qué es un videojuego de rol necesitaremos primero entender \textit{qué es un videojuego} y \textit{qué es un juego de rol}.

\medskip

Definir qué es un videojuego es una tarea bastante complicada, sobre todo teniendo en cuenta los matices con los que podemos definirlo (académicos, de diseño, experimentales o tecnológicos). Una de las definiciones más acertadas nos la da \cite{EspositoVJ}, que lo define como \comillas{un \textit{juego} que se \textit{juega} gracias a un \textit{aparato audiovisual}, y que puede estar basado en una \textit{historia}}. El propio \citeauthor{EspositoVJ} se encarga de definir qué es el \textit{juego}, qué es \textit{jugar}, qué es el \textit{aparato audiovisual} y qué es la \textit{historia}. La diferencia fundamental de los juegos tradicionales frente a los videojuegos es la existencia de ese \textit{aparato audiovisual} (videoconsolas, ordenador, dispositivos móviles) que pueda ofrecer una interacción \comillas{humano-máquina}, ya que es esta interacción recíproca la que hace que los videojuegos sean videojuegos y no otro tipo de entretenimiento multimedia.

\medskip

Los juegos de rol, al haber nacido desde juegos de mesa tradicionales, han tenido el suficiente tiempo como para poder desarrollarse y poder ajustarse a una definición mucho más precisa, aunque, nuevamente, debido a la amplia variedad de juegos que se pueden catalogar como \comillas{de rol}, es posible que hasta la definición más completa no pueda abarcarlos a todos. \citeauthor{LortzRPG} (\citeyear{LortzRPG}, citado en \cite{FineRPG}), define los juegos de rol como \comillas{todos aquellos juegos que permiten a un determinado número de jugadores asumir los roles de personajes imaginarios y operar con cierto grado de libertad en un entorno imaginario}.

\smallskip

Por otra parte, \citeauthor{TychsenRPG} (\citeyear{TychsenRPG}, citado en \citet*{RPGH&D}) enumera los elementos imprescindibles que un juego debe tener para ser considerado como juego de rol:
\begin{itemize}
	\item Narrar una historia adecuándose a una serie de reglas. Tanto la historia como las reglas son únicas para cada juego.
	\item Reglas, múltiples participantes (al menos dos) y un mundo ficticio sobre el que se va a desarrollar la acción. Todos los participantes deben conocer previamente la ambientación, el entorno y las reglas.
	\item La gran mayoría de participantes controlarán, como mínimo, a un personaje durante toda la partida, y será con ese o esos personajes con quienes interactuará con el entorno.
	\item La existencia, por lo general, de un \textit{director de juego} (GM, \textit{gamemaster}), que será el responsable de gestionar aquellos elementos del juego o del entorno ficticio que no se encuentran bajo el control directo de los jugadores.
\end{itemize}

\medskip

Ahora que ya sabemos \textit{qué es un videojuego} y \textit{qué es un juego de rol}, y pese a que no podemos dar una definición válida para todos los videojuegos que estén en esta categoría, definiremos un videojuego de rol (RPG, \textit{role-playing game}, o también, más raramente, CRPG, \textit{computer role-playing game}) como todo aquel programa \textit{software}, con carácter lúdico, que permita a sus usuarios encarnar el rol de uno o varios personajes en un mundo ficticio, donde pueden tomar decisiones, interactuar con su entorno o mundo con cierta libertad y desarrollar una narrativa para conseguir un determinado objetivo.

\smallskip

Como se ha mencionado anteriormente, esta definición vale para la gran mayoría de RPG; pero, si pensamos en los videojuegos que se venden cada día en las tiendas, una proporción significativa de estos puede encajar dentro de la definición anterior sin necesariamente tener que ser un RPG, ya que prácticamente en todos se encarna el rol de un personaje, casi todos permiten una interacción con el entorno y todos tienen un objetivo final. Esto nos lleva a plantearnos si las clasificaciones de videojuegos según su género son demasiado limitadas para los videojuegos actuales.

\subsection{El problema de los géneros de videojuegos}
Antes de adentrarnos en los RPG, veremos por qué las distintas clasificaciones de videojuegos atendiendo a características similares tienen bastantes lagunas a la hora de categorizar las entregas más modernas.

\medskip

Los primeros intentos de clasificar los videojuegos mediante características comunes surgen en la década de los 80, principalmente por diseñadores y desarrolladores, como \cite{Crawford84}, quien los categorizó entre aquellos que \textit{requieren de habilidades previas por parte del usuario} (como juegos de combate o de carreras), y \textit{juegos de estrategia} (que engloba al resto de juegos, como los de aventuras, los educativos, y los RPG). Estas categorías son \textit{funcionales}, se centran en las mecánicas de juego y enfatizan cómo los jugadores interactúan con el sistema.

\smallskip

En la actualidad, los distintos géneros vienen dados por una mezcla de mecánicas, temas, elementos narrativos, estética, lugar de origen y plataforma, y están gravemente influenciados por las tendencias del mercado, discursos mediáticos y la propia percepción de los jugadores. También, muchas veces se quiere categorizar a los videojuegos con etiquetas propias de otras modalidades, como el cine o la literatura, que son incapaces de capturar los aspectos únicos que definen a un videojuego.

\smallskip

\cite{FailGeneros} argumentan que las categorías existentes a día de hoy se quedan cortas para satisfacer los propósitos del género en videojuegos (identidad, agrupación, \textit{marketing} y educación), ya que dada la gran diversidad de juegos que tenemos ahora, resulta casi imposible poder agrupar la naturaleza multifacética de estos en etiquetas tan \comillas{tradicionales}.

\smallskip

Pongamos el ejemplo de uno de los juegos más exitosos de la última década, \citegame{The Legend of Zelda: Breath of the Wild}{BOTW}, que mezcla elementos de libre exploración (lo que lo convertiría en un \textit{sandbox} o \comillas{juego libre}), acción en tiempo real e investigación y resolución de rompecabezas (lo que lo convertiría en un juego de \comillas{acción-aventura}), y la progresión en tiempo real del personaje característica de los RPG (puedes ir consiguiendo nuevas habilidades o mejorando el equipamiento). Es por eso que muchas veces tenemos géneros intermedios, como en este caso, que podríamos considerar a \textit{Breath of the Wild} como un RPG de acción (ARPG, \textit{action role-playing game}, que igualmente siguen sin englobar a la increíble diversidad de juegos. Este problema también se puede aplicar a la inversa, donde juegos como \citegame{Undertale}{Undertale}, esencialmente un RPG, tiene elementos, como el combate, propios de otro género de juegos.

\smallskip

Hay muchas formas de evitar este problema, y quizá la mejor solución sea dejar de considerar a los géneros como \comillas{cajones estancos} en los que un videojuego no pueda pertenecer a dos de estas categorías simultáneamente \citep{Apperley}, sino que sean más bien un espectro, sin fronteras establecidas, en el que un videojuego pueda caer entre dos categorías distintas.

\smallskip

En resumen, antes de categorizar un videojuego en según qué género, hay que entender que resulta imposible que este se reduzca a una fórmula fija, sino que hay que entenderlo como una combinación fluida de mecánicas, elementos narrativos, temas y estética, que varían de juego a juego.

\subsection{Historia de los videojuegos de rol}
Es a mediados de la década de los 70 cuando podemos hablar del nacimiento de los videojuegos de rol. Dadas las limitaciones tecnológicas de la época, los RPG primitivos no eran más que simples adaptaciones de juegos de mesa ya existentes por entonces, como \citegame{dnd}{dnd}, una adaptación de \citegame{Dungeons \& Dragons}{ogdnd}, quizás el juego de rol más famoso, que mezcla mecánicas de combate con temas fantásticos. De \textit{Dungeons \& Dragons} se ha adoptado en la gran mayoría de juegos el lanzar los dados, las estadísticas de los personajes (por ejemplo, inteligencia o fuerza), o sistemas de niveles. 

\smallskip

Estos primeros juegos contaban con interfaces basadas en texto, y normalmente estaban programados en los grandes ordenadores que se encontraban en campus universitarios como los de Harvard o la Universidad de Illinois. Esta primera etapa, \cite{barton2008dungeons} la denomina como la etapa de \comillas{años oscuros} (en un símil con los años oscuros del Medievo europeo), principalmente por lo anteriormente mencionado, pero también por la escasa información que tenemos hoy en día sobre estos juegos, ya que muchos no se han conservado son \textit{lost media} (aquellos materiales multimedia que ya no existen en ningún formato y para los cuales no hay ninguna copia disponible).

\smallskip

De esta primera etapa también cabe mencionar tres videojuegos que dieron comienzo a tres distintos subgéneros dentro de los RPG: 
\begin{itemize}
\item \citegame{Rogue}{rogue}, que dio lugar a los videojuegos de mazmorra o \textit{roguelikes}, caracterizados por la generación aleatoria de un laberinto o mazmorra en el que se desarrolla una aventura basada por turnos. En este tipo de juegos la muerte es permanente, y al perder la partida se empieza en una nueva desde cero. 
\item \citegame{Wizardry}{wizardry}, que dio lugar a los videojuegos de exploración de mazmorra o \textit{dungeon crawlers}, similares a los anteriores, pero centrados en la exploración de la mazmorra, con un énfasis en la progresión de los personajes y de la gestión de la \textit{party} o escuadrón (el conjunto de personajes que juntos intentan alcanzar objetivos comunes). 
\item \citegame{Ultima}{ultima}, que dio lugar a los RPG de mundo abierto. En esta clase de juegos, los jugadores pueden explorar libremente ciudades, mazmorras, bosques y otro tipo de entornos, manteniendo las mecánicas tradicionales de otros RPG. 
\end{itemize}

\medskip

Con el salto tecnológico que hubo a mediados de la década de los 80, los RPG comienzan a separase cada vez más de ser meras adaptaciones de juegos ya existentes a desarrollar sus propias historias. A partir de esta época encontramos dos corrientes bastante diferenciadas de RPG, los \comillas{occidentales} (WRPG, \textit{Western role-playing game}), con más libertad de decisión para los jugadores tanto en personalización como en la historia, y con temáticas realistas (como los anteriormente mencionados \textit{Rogue}, \textit{Wizardry} y \textit{Ultima}); y los \comillas{orientales} o \comillas{japoneses} (JRPG, \textit{Japanese role-playing game}, por ser Japón el país que más videojuegos de este tipo produce), centrados en una narrativa lineal con temáticas fantásticas y mecánicas basadas por turnos. Dos grandes videojuegos que definieron el género de los JRPG son \citegame{Dragon Quest}{dq} y \citegame{Final Fantasy}{ff}, cuyas sagas permanecen hasta nuestros días con nuevas entregas cada pocos años. Son estos años de apogeo de los RPG los que \citeauthor{barton2008dungeons} denomina como \comillas{etapa dorada}.

\medskip

La década de los 90 supuso un grave declive para los RPG, especialmente en los mercados occidentales, ya que la aparición de juegos de acción en 3D, como \citegame{Doom}{doom}, \citegame{Quake}{quake} o \citegame{Tomb Raider}{lc}, hizo que el mercado cambiase hacia este tipo de juegos, mucho más rápidos y frenéticos que los RPG, que se consideraban obsoletos, con una pesada carga textual y mucho más lentos de jugar. También, la aparición de videoconsolas mucho más potentes y baratas, como la \textit{PlayStation} (Sony, 1994) o la \textit{Nintendo 64} (Nintendo, 1996), para las cuales los RPG occidentales no tenían portabilidad, y los altos costes de desarrollo y producción que supuso el cambio de cartuchos tradicionales a nuevos formatos como CD-ROM, hicieron que los RPG, especialmente los WRPG, tuviesen este gran declive que solo pudo recuperarse con la entrada del nuevo milenio.

\medskip

En el nuevo milenio, que para \citeauthor{barton2008dungeons} es la \comillas{etapa de platino}, surgen sagas con juegos con gráficos mucho más sofisticados y con narrativas mucho más profundas, como la saga \textit{The Elder Scrolls}, más concretamente, su tercera entrega, \citegame{Morrowind}{tes}; la saga \citegame{Fallout}{fallout}; la saga \citegame{Baldur's Gate}{baldurs}; o la saga \citegame{Diablo}{diablo}, que llevan hasta el límite las propias definiciones del género por las mezclas con otros géneros (llegando a ser juegos \comillas{híbridos}). La tendencia de los mundos abiertos continúa y se mejora, llegando a ser algunos RPG como \citegame{The Elder Scrolls V: Skyrim}{skyrim} o \citegame{The Witcher 3: Wild Hunt}{witcher} algunos de los juegos más jugados de la historia, e incluso ganando múltiples premios a juego del año.

\section{Sobre el desarrollo de videojuegos}

\subsection{Motor de videojuegos}

\subsubsection{Componentes de un motor de videojuegos}

\subsubsection{Separación entre motor y \textit{gameplay}}

\subsubsection{Programación dirigida por datos (\textit{DDP}) en videojuegos}

\subsubsection{Programación multiplataforma}

\subsection{Editor de videojuegos}

\subsubsection{Editores específicos para desarrollo de \textit{RPG}}


\chapter{Descripción del Trabajo}
\label{cap:descripcionTrabajo}

Aquí comienza la descripción del trabajo realizado. Se deben incluir tantos capítulos como sea necesario para describir de la manera más completa posible el trabajo que se ha llevado a cabo. Como muestra la figura \ref{fig:sampleImage}, está todo por hacer.

\begin{figure}[h]
	\centering
	\includegraphics[width = 0.5\textwidth]{Imagenes/Vectorial/Todo.pdf}
	\caption{Ejemplo de imagen}
	\label{fig:sampleImage}
\end{figure}

Si te sirve de utilidad,  puedes incluir tablas para mostrar resultados, tal como se ve en la tabla \ref{tab:sampleTable}.


\begin{table}
	\centering
	\begin{tabular}{c|c|c}
		\textbf{Col 1} & \textbf{Col 2} & \textbf{Col 3} \\
		\hline\hline
		3 & 3.01 & 3.50\\
		6 & 2.12 & 4.40\\
		1 & 3.79 & 5.00\\
		2 & 4.88 & 5.30\\
		4 & 3.50 & 2.90\\
		5 & 7.40 & 4.70\\
		\hline
	\end{tabular}
	\caption{Tabla de ejemplo}
	\label{tab:sampleTable}
\end{table}

%\include{Capitulos/Capitulo4}
%\include{Capitulos/Capitulo5}
\chapter{Evaluación y Conclusiones}
\label{cap:conclusiones}

Al conocer en profundidad cómo funcionan todos los sistemas del motor y el editor, no podemos tener la certeza de que estos funcionen como un usuario promedio esperaría. Por esto, realizar pruebas con usuarios sin sesgos es crucial para obtener la información de las posibles brechas entre los programas y esos usuarios. De cara a ofrecer una experiencia de uso agradable se analizarán las pruebas en busca de formas de mejorarla. Se buscarán distintos tipos de perfiles para obtener una mayor cantidad de datos sobre los posibles conflictos que estos tengan con los sistemas.

\section{Objetivo de la evaluación}
Los objetivos de las pruebas con usuarios son principalmente dos. El primero es evaluar si los usuarios entienden los sistemas y son capaces de usarlos cómodamente, es importante que el entorno sea agradable de usar para permitir que los usuarios entiendan las capacidades del mismo y las puedan aprovechar. El segundo es evaluar si cada usuario es capaz de usar el editor creativamente, de forma que los videojuegos que cree un usuario suponga una experiencia diferente a los que haga otro, comprobando así la flexibilidad de uso del editor.

\section{Metodología}
Para realizar estas pruebas, una vez se plantearon los objetivos, se diseñó cómo se iban a realizar. Cada usuario realizará las pruebas de manera individual en compañía de un supervisor, tendrán una duración estimada de una hora. A cada probador se le preguntan lo primero de todo su experiencia tanto jugando como creando videojuegos, haciendo énfasis en los RPG. Una vez con estos datos se le entregan varios ficheros al empezar la prueba, el ejecutable del editor, un paquete de recursos para no desaprovechar tiempo de la prueba buscando y una breve guía de uso del editor. En esta guía se explica brevemente el modo de uso de cada uno de los sistemas disponibles en el editor. Una vez con esto, el probador seguirá la guía para entender los mecanismos y podrá para utilizarlos a su libre albedrío.

\medskip

El investigador se encargará de registrar todas las interacciones destacables de los probadores con el editor, además de prestar ayuda si surge algún error con el programa a lo largo de las pruebas. Para poder observar una experiencia real, el investigador no podrá intervenir en las acciones del usuario más que en esos casos excepcionales con los errores o en momentos de bloqueo en los que el probador no es capaz de continuar por alguna brecha. Estas pruebas están pensadas para realizarse tanto en persona como telemáticamente a mayor comodidad de cada caso concreto. Una vez finalice la prueba se le pedirá al usuario una valoración general sobre los sistemas y su experiencia de uso aportando sus quejas e ideas.

\section{Resultados}

Los usuarios pueden agruparse en 3 categorías según su experiencia con los videojuegos y, en particular, con los RPG.

\smallskip 

En primer lugar, se encuentran los usuarios más ajenos al mundo de los videojuegos, con poca experiencia jugando a juegos en general y a juegos RPG en concreto. Con este tipo de usuario el objetivo es ver como de intuitivas son las herramientas y la UI de la aplicación del editor, así como ver si se entienden bien las instrucciones de las pruebas. Este primer grupo consiguió crear mapas con cierto sentido, pero la mayoría utilizó únicamente una o dos capas por lo que con esto quedaban un poco limitados a la hora del diseño. No invirtieron gran cantidad de tiempo explorando los tilesets e hicieron unos diseños más bien sencillos. Las configuraciones las entendieron bastante bien, y siguiendo la guía las consiguieron rellenar. A la hora de crear eventos se quedaron en los behaviours más sencillos y no intentaron crear eventos complejos usando variables y condiciones. En general la prueba fue satisfactoria y los resultados estaban dentro de lo esperado. 

\medskip

En segundo lugar, se tienen a usuarios que sí juegan a videojuegos y algunos incluso están relacionados con la industria, pero no tienen conocimientos de programación o de desarrollo más allá de sus áreas específicas. Estos si aprovecharon mucho más el apartado del mapeado, utilizando más capas y varios tilesets creando mapas mucho más interesantes. A la hora de los eventos se les insistió más que al primer grupo para que intentaran crear algunos eventos algo más complejos, pero al parecer no les quedaba tan claro como estructurarlos para aprovechar al máximo las funcionalidades y la mayoría se quedó en cosas algo más sencillas o tuvieron que recibir bastante ayuda por parte de los supervisores. Aunque los eventos no requieren conocimientos de programación, su estructura es similar a la de un sistema basado en lógica condicional, lo que exige más tiempo del que se disponía durante la prueba para entender cómo organizarlos correctamente. 

\medskip

Por último, está el grupo de gente que tiene conocimiento en el desarrollo de los videojuegos, incluso en el desarrollo de \textit{RPGMaker} específicamente, lo que fue muy interesante a la hora de obtener \textit{feedback} comparativo. Este grupo entendió bien las herramientas de mapeado y señaló que echaba en falta algunas más avanzadas que agilizaran un poco el proceso, como la opción de seleccionar varios \textit{tiles} a la vez o de rellenar zonas con \textit{tiles} de un tipo. También señalaron lo lento que es el proceso de creación de \textit{sprites} y animaciones, proponiendo algún tipo de automatización de este como ocurre con los \textit{tilesets}. En cuanto a los eventos lograron entender como montar estructuras más avanzadas, lo cual fue un resultado positivo. 

\medskip

Adicionalmente a todas las observaciones y el \textit{feedback} recibido, las pruebas también fueron útiles para detectar varios \textit{bugs} y cosas que no estaban funcionando correctamente, lo que fue de gran ayuda de cara a dar un pulido final a la versión del editor. 

\section{Análisis de los resultados y conclusiones}
Aunque las pruebas no alcanzaron la extensión ni la variedad deseables debido a limitaciones logísticas, se pueden sacar algunas conclusiones interesantes de cara al desarrollo actual y sobre todo a un posible trabajo futuro.

\smallskip 

Una de las prioridades identificadas es la mejora de la experiencia de diseño de mapas y de creación de \textit{sprites} y animaciones. Con este objetivo, se proponen las siguientes mejoras funcionales:

\begin{itemize}
	\item Incorporar la posibilidad de seleccionar varios \textit{tiles} simultáneamente y pintarlos en el mapa de forma conjunta. 

	\item Añadir una herramienta para rellenar automáticamente una zona completa con el \textit{tile} seleccionado. 

	\item Permitir la selección de \textit{tiles} directamente desde el mapa, sin necesidad de localizarlos manualmente en el \textit{tileset}. 

	\item Incluir una función para deshacer la última edición realizada. 

	\item Automatizar el proceso de creación de múltiples \textit{sprites} y su conversión en animaciones a partir de una única fuente, de manera similar al funcionamiento del editor de \textit{tilesets}. 
\end{itemize}

Además, se considera conveniente implementar mejoras en forma de \textit{tooltips} en la sección de eventos, con el objetivo de explicar de manera clara todas las funcionalidades disponibles. También se sugiere incluir eventos predefinidos que sirvan como ejemplos prácticos para facilitar su comprensión y uso.




%%%%%%%%%%%%%%%%%%%%%%%%%%%%%%%%%%%%%%%%%%%%%%%%%%%%%%%%%%%%%%%%%%%%%%%%%%%
% Si el TFG se escribe en inglés, comentar las siguientes líneas 
% porque no es necesario incluir nuevamente las Conclusiones en inglés
\begin{otherlanguage}{english}
\chapter*{Introduction}
\label{cap:introduction}
\addcontentsline{toc}{chapter}{Introduction}

\chapterquote{I had this really bizarre conversation once with a couple of lawyers and they were talking about \comillas{How do you pick your target market? Do you use focus groups and poll people and all this?} It's like \comillas{No, we just write games that we think are cool.}}{John Carmack}


\section*{Motivation}
Role-playing video games are one of the most influential genres in the industry\footnote{In total, more than one billion copies of role-playing video games have been sold throughout history, according to \textit{VGChartz} (\url{https://shorturl.at/VI2Zy}).}, from their origins in the 1980s to the present day, where they account for a large portion of the market share\footnote{According to \textit{Rocket Brush Studio}, role-playing video games correspond to the best-selling genre in the mobile market, the third-best in the PC market, and the fifth-best in the console market (\url{https://rocketbrush.com/blog/most-popular-video-game-genres-in-2024-revenue-statistics-genres-overview}).}.

\smallskip

The development of this type of game has undergone significant changes over the years, along with consequent improvements in hardware and software, which have allowed us to evolve from simple monochrome scenes with text-based interfaces to real-time 3D graphics, complex open worlds, and advanced interaction systems. These improvements have influenced not only the visual aspect but also the narrative, gameplay, and player immersion, enabling increasingly rich and personalized experiences.

\medskip

Numerous video game development and editing tools have appeared in the last two decades, but the more commercial ones are designed for games of all kinds. As a result, they tend to offer more general features that are not closely related to role-playing game development, and they often restrict certain functionality, making the development process more complicated.

\smallskip

Those tools that are specifically designed for role-playing game development have two fundamental flaws:
\begin{itemize}
	\item Those that offer an intuitive, easy-to-use interface that is friendly to newcomers in game development are locked behind a paywall which, although not very expensive, forces amateur users to pay for a software they will rarely unless they consolidate as developers.
	\item Those that are not locked behind a paywall or are open source, often have interfaces and systems that are difficult for beginners to understand, leading many to abandon game development due to the complexity of these tools.
\end{itemize}

\medskip

The main motivation behind this Project arises from the need to have easy-to-use and flexible tools, both for experimented independent developers who desire to create games without the limitations imposed by general-purpose engines, and for people with little knowledge of video game development or programming.

\section*{Objectives}
This project aims to develop a 2D game engine specifically designed for role-playing games, along with an editor that enables simple and straightforward game development for this engine. The editor will be capable of generating executable files that users can distribute without the need for any additional steps after development.

\smallskip

The editor's interface will be designed for beginners in game development, while still alowing more experienced users to create larger and more ambitious projects.

\medskip

To accomplish this:
\begin{itemize}
	\item A cross-platform engine that allows the user to implement role-playing games will be developed.
	\item A cross-platform editor that eases the implementation of a role-playing game in the developed engine will be developed.
	\item The Project will be tested with users to demonstrate its functionality, and necessary conclusions will be drawn, along with potential improvements for the future.
\end{itemize}

\section*{Work Plan}
\begin{figure}[t]
	\scalebox{0.9}{
		\begin{ganttchart}[
    		hgrid,
    		vgrid,
    		time slot format=isodate-yearmonth,
    		title label font=\small,
    		time slot unit=month,
    		title/.style={fill=blue!20},
    		title height=1.2,
    		x unit=2cm
  			]{2025-01}{2025-05}
	
			\gantttitle{January}{1}
			\gantttitle{February}{1}
			\gantttitle{March}{1}
			\gantttitle{April}{1}
			\gantttitle{May}{1}
			\gantttitle{June}{1}\\
	
			\ganttbar{Research on the current state}{2025-01}{2025-02}\\
			\ganttbar{Project design}{2025-01}{2025-02}\\
			\ganttbar{Engine development}{2025-02}{2025-05}\\
			\ganttbar{Editor development}{2025-02}{2025-05}\\
			\ganttbar{Document writing}{2025-03}{2025-05}\\
			\ganttbar{User testing}{2025-05}{2025-05}
		\end{ganttchart}
	}
	\caption{Gantt diagram showing the temporal planning of the work.} 
	\label{fig:ganttingles}
\end{figure}

To accomplish the previous objectives, the work plan will be divided in three phases: 
\begin{itemize}
	\item Research on the current state of engines and editors specifically designed for role-playing games, as well as on role-playing games themselves. This section will aim to identify the common features shared by all engines, editors and games, both open-source and commercial ones, and propose improvements to address the issues these tools may present for inexperienced new users. 
	\item Project Design. Based on the common features identified in the previous phase, an initial design will be proposed to serve as a basis for the development of the Project. This design, while not immutable, should be as close to final as possible to avoid any issues during the development phase.
	\item Project Development. Once the design phase is completed, development will begin. This phase will, in turn, be split across multiple phases:
		\begin{itemize}
			\item Setting up the development environment. A development environment will be chosen based on the Project's needs, and everything will be configured to minimize effort during the development of the applications.
			\item Engine development. An engine will be developed according to the previously established plans, with cross-platform support for both PCs and Android mobile devices.
			\item Editor development. Like the engine, the editor will be developed according to the established design. It must support multiple platforms on PC, namely Windows, MacOS and Linux, but not Android.
		\end{itemize}
	\item User testing. To ensure the correct functioning of the Project, the final tools will be tested by a variety of users not involved in their development. After the tests are completed, the necessary conclusions will be drawn, and the Project will be adjusted to address any critical issues that require attention before releasing the public version. Features that would require excessive effort to implement within the available time will be left as future work.
\end{itemize}


\section*{Tools and Methodology}
Regarding the tools, Git will be used as the version control system, with a repository hosted on GitHub and managed through GitHub Desktop. Task management will be carried out using GitHub Projects, available on the GitHub website.

\smallskip

Access to the repository containing the code may be done through this link: \href{https://github.com/almasso/rpgbaker}{\textbf{https://github.com/almasso/rpgbaker}}.

\medskip

Regarding the development tools for the Project, CLion will be used as the C++ development environment, and Android Studio for Android development in Java. CMake will be used as the tool for generating external libraries, and MinGW for handling compilation on Windows, Clang on MacOS, and GCC on Linux. Detailed reasons for choosing these tools can be found in section \ref{sec:decisiones}.

\medskip

On the other hand, PDF generation of the document will be carried out using \LaTeX, with the \texis template, and Texmaker will be used as the editing environment.

\bigskip

Regarding the working methodology, bi-weekly meetings will be proposed with the tutor, although the number of weeks may vary depending on the progress made. During these meetings, the current status of the Project will be presented, and advice related to the design or development of the Project will be requested.

\medskip

Communication with the team will be established through both in-person meetings when important issues need to be addressed and via software that allows messaging and voice chat, such as Discord. This tool will also be used for conducting user testing.

\section*{Document Structure} 
In Chapter \ref{cap:juegosrol}, \textit{Role-Playing Games}, the topic of role-playing games is discussed, including their characteristics, origins, and an analysis of some of the most important titles in the genre.

\medskip

In Chapter \ref{cap:videojuegosrol}, \textit{Role-Playing Video Games}, role-playing video games and their history are introduced, along with the development process and the use of game engines and editors.

\medskip

In Chapter \ref{cap:planteamiento}, \textit{Project Planning}, the design and decisions made prior to the development of the applications are discussed in detail.

\medskip

In Chapter \ref{cap:desarrollo}, \textit{Project Development}, the internal structure of the engine and editor in terms of code are discussed, while addressing the problems encountered during this phase and the solutions proposed.

\medskip

In Chapter \ref{cap:conclusiones}, \textit{Evaluation and Conclusions}, the research questions and objectives are presented, along with the development of the user testing phase and an analysis of the tests, from which conclusions are drawn.

\medskip

Finally, in Chapter \ref{cap:trabajoFuturo}, \textit{Future Work}, potential improvements for the final implementation of the applications that may be of interest are detailed.









\chapter*{Conclusions and Future Work}
\label{cap:conclusions}
\addcontentsline{toc}{chapter}{Conclusions and Future Work}

Conclusions and future lines of work. This chapter contains the translation of Chapter \ref{cap:conclusiones}.



\end{otherlanguage}
%%%%%%%%%%%%%%%%%%%%%%%%%%%%%%%%%%%%%%%%%%%%%%%%%%%%%%%%%%%%%%%%%%%%%%%%%%%

\chapter{Contribuciones Personales}
\label{cap:contribucionesPersonales}
\addcontentsline{toc}{chapter}{Contribuciones Personales}

\section*{Miguel Curros García}
Mis tareas en el proyecto se centraron en el diseño y desarrollo del motor genérico de videojuegos así como dentro del apartado de \textit{gameplay} todo lo relacionado con el sistema de eventos, tanto a nivel de motor como de editor. Adicionalmente, también creé la gestión de persistencia con lectura y escritura de los archivos de datos específicos del editor y parte del proceso de \textit{build}.

\medskip

En primer lugar, participé en conjunto con Alejandro González en la fase de diseño del motor donde dejamos bien claras cuáles serían las distintas capas de abstracción de la implementación. Desde el principio teníamos claro que la base del motor iba a estar basada en una arquitectura EC (\textit{Entity-Component}), construyendo el resto de las funcionalidades alrededor. Comenzamos a definir cómo estructuraríamos esa arquitectura a través de un diagrama \textit{UML} donde definimos tanto los distintos módulos que formarían el motor como la forma en la que se conectarían las partes más básicas de este con las funcionalidades específicas de \textit{gameplay}.

\medskip
 
A continuación, comencé con el desarrollo de este centrándome en los apartados de carga y gestión de recursos, audio y colisiones. Comenzando con la gestión de recursos, esta es una parte crucial del motor al estar este dirigido por datos; todo aquello que se quiere que ocurra en un juego estará definido en los datos a interpretar al motor.
Esta parte tiene 3 partes clave: \texttt{Resource}, \texttt{ResourceHandler} y \texttt{ResourceMemoryManager}. Cada una de esas clases se encarga de una parte clave de la carga y descarga de recursos. 
\begin{itemize}
	\item \texttt{Resource} servirá de clase base para todos los tipos de recursos que se quieran cargar, cada uno de ellos implementará cómo se carga y descarga desde los archivos de datos.
	\item \texttt{ResourceHandler} será a quien recurran las distintas partes del motor que necesiten cualquier tipo de recurso. A través de la ruta a un recurso de un tipo especificado dará acceso a una instancia conteniendo ese mismo recurso cargado.
	\item Por último, \texttt{ResourceMemoryManager} se encargará de gestionar cuánta memoria está siendo ocupada en este momento por los recursos. A partir de un límite especificado en un archivo de configuración esta clase se encargará de que no se supere ese umbral de memoria máxima ocupada por los recursos. Si se solicitara un recurso y no hubiera memoria suficiente para cargarlo a través de un algoritmo LRU se irá liberando memoria de otros recursos hasta que haya suficiente para cargar el nuevo. Con esto, conseguí estructurar la carga de recursos bajo demanda que evitaría posibles largos tiempos de espera los juegos.
\end{itemize}

\medskip

Una vez esto estaba listo desarrollé el sistema de sonido que permitiría el control sobre la reproducción de archivos de audio dentro de los juegos. Para esto decidimos usar el nuevo sistema de audio de \textit{SDL3} que cumplía todas nuestras necesidades. Nuestro objetivo con este sistema era la implementación de un componente \texttt{AudioSource} a través del que gestionar la reproducción de un archivo de audio. Además, queríamos que cada uno de estos \texttt{AudioSource} pudieran estar asociados a un conjunto de sonidos, pudiendo tener de esta forma un control más profundo del volumen; pues es típico y útil en los videojuegos permitir al jugador controlar el volumen general, de la música o de los efectos por separado.

\smallskip

Para poder reproducir un sonido creé una clase \texttt{AudioClip} que se encargaría de envolver las funcionalidades de \texttt{SDL\_AudioStream}; la base de la reproducción de sonidos en \textit{SDL3}; y ofrecer una interfaz acorde a nuestras necesidades. Para poder asignar estas pistas a distintos conjuntos de sonidos implementé la clase \texttt{AudioMixer}, que tendría un control de volumen asociado que se aplicaría a cada uno de sus \texttt{AudioClip}. Una vez con estas clases básicas pude implementar \texttt{AudioSource} que las utilizaría para ofrecer su funcionalidad esperada.

\medskip

Por último en la base del motor, para la detección y gestión de colisiones buscábamos algo sencillo. Por el tipo de videojuegos que ofrecemos implementar basta con comprobaciones de colisiones entre formas rectangulares. Son comprobaciones sencillas y \textit{SDL} ya ofrece funciones para comprobar si dos rectángulos tienen intersección. Con esto, la implementación de un componente \texttt{Collider} fue sencilla. A través de una clase \texttt{CollisionManager}, en cada actualización del juego se comprobarían las colisiones entre estos \texttt{Collider}. En estas comprobaciones se guardaría la información de qué \texttt{Collider} está colisionando, acaba de empezar a colisionar o acaba de dejar de colisionar con otro. De esta manera se ofrece un control sencillo para implementar comprobaciones dependientes de colisiones entre elementos de juego.

\medskip

Una vez terminamos el apartado del motor genérico comenzamos con el desarrollo de partes específicas de \textit{GamePlay} donde centré mis aportaciones en el sistema de eventos. De cara a ofrecer una experiencia personalizable de crear un videojuego sin necesidad de programar un punto clave era nuestro sistema de eventos. Este se centra alrededor del componente \texttt{EventHandler} que sirve como conjunto de eventos de una entidad. Cada \texttt{Event} está compuesto por una condición, \texttt{EventCondition}, y comportamientos, \texttt{EventBehaviour}.
\begin{itemize}
	\item Comenzando por las condiciones, estas se encargan de manifestar si el estado del juego es el que esperan. Para ello, cada una de las condiciones debe implementar sus propias comprobaciones, esto lo hice a través de un sistema de herencia. Además, igual que pasa con los componentes, cada una de estas condiciones deben poder crearse a partir de los datos proporcionados por el juego, por esta razón creé una clase \texttt{EventConditionFactory}, que se encargaría de instanciar el tipo correspondiente de \texttt{EventCondition} dado un identificador.
	\item Los comportamientos decidimos que estuvieran implementados en \textit{Lua}. De esta manera se podrían ampliar las funcionalidades de un juego sin necesidad de recompilar el motor, ofreciendo una experiencia más flexible. Para implementarlo de la manera que diseñamos era necesario poder tener algún tipo de POO (Programación orientada a objetos) en \textit{Lua}, pero no es algo que ofrezca el lenguaje por defecto. A través de asignar funciones a tablas y modificando sus \textit{metatablas} pude obtener un comportamiento similar a las clases y la herencia. A partir de esa base implementé la clase \texttt{EventBehaviour} en \textit{C++} envolviendo a cada instancia que hubiera en una escena de los distintos comportamientos implementados en \textit{Lua}. Además de esto, para poder implementar cada uno de los comportamientos fue necesario definir desde \textit{C++} qué clases y cómo se podían modificar dentro de las implementaciones de estos en \textit{Lua}.
\end{itemize}

\medskip

Cuando esto estuvo listo comencé con el desarrollo del editor donde creé la gestión de persistencia con lectura y escritura de los archivos de datos específicos del editor. Estos datos decidimos que íbamos a guardarlos también en \textit{Lua} pues nos sería más sencillo de implementar al tener ya la infraestructura montada para ello. Gracias a la clase del editor \texttt{LuaManager} creada por Alejandro Massó, que permite acceder a tablas de \textit{Lua} de un archivo, crear nuevas y guardar estas en nuevos archivos, la tarea consistió en escribir y leer estas tablas. Cada vez que se quisiera salvar un proyecto cada uno de sus recursos guardarían sus parámetros en nuevas tablas de \textit{Lua} que se escribirían en archivos en subdirectorios específicos dentro de un directorio \comillas{projectfiles} dentro de la ruta del proyecto.

\medskip
 
A continuación desarrollé todo el apartado de Eventos, desde la creación, edición y posterior traducción a archivos preparados para ser leídos por el motor. Para este apartado lo primero que hice fue crear una nueva pestaña en el editor, la pestaña de edición de eventos. En esta se podría ver un desplegable donde escoger un evento, una sección donde se mostraría la condición del evento y otra sección donde se mostrarían los comportamientos. 
La parte de creación y selección de eventos fue sencilla, aprovechando que Alejandro Massó ya había implementado la creación y selección de mapas, reutilicé el código que pude para esta parte. Para hacer los apartados de la condición y los comportamientos, igual que con la condición en el motor, tuve que crear factorías que permitieran crear los distintos tipos de instancias de cada uno de estos a partir de identificadores. Esto es porque cada condición y cada comportamiento debe definir en el editor su propia clase, ya que su persistencia e interfaz gráfica difieren entre sí, de modo que requieren implementaciones distintas; además por esa misma persistencia es por lo que se necesita la factoría, cada evento puede tener cualquier tipo de condición o comportamiento y se debe poder reconstruir al abrir un proyecto guardado.

\medskip

Una vez completé todas las interfaces gráficas y la persistencia de todos los tipos de condiciones y comportamientos comencé con el proceso de \textit{build}; convertir estos datos a la sintaxis que el motor reconoce. En el caso de las condiciones fue sencillo, pues el formato en el que guardan su información en el editor es casi idéntico a lo que necesita el motor. Lo que refiere a los comportamientos no es así, hubo dos puntos clave en este apartado: dependencias de componentes y formato de escritura. Algunos de los comportamientos en el motor funcionan por su cuenta, es decir, con que existan dentro de su \texttt{EventHandler} cumplirán su función sin problema; en cambio hay otros que necesitan otros componentes para funcionar también de modo que, para que su proceso de \textit{build} fuera correcto necesitaban indicarle a su entidad que escribiera los componentes faltantes con los parámetros correctos. Por último, a diferencia del resto de proceso de \textit{build} este apartado no podía hacer uso del mecanismo que usamos para escribir el resto de tablas de \textit{Lua}. Esto ocurre porque no estamos añadiendo otra tabla al uso, se necesita la llamada a la construcción de ese \texttt{EventBehaviour} porque el mecanismo de persistencia que se usa en el resto del editor no graba las \textit{metatablas} asociadas a cada tabla, y esto es crucial de cara al funcionamiento de los comportamientos.

\section*{Alejandro González Sánchez}
Mi contribución al Proyecto se ha centrado en el diseño e implementación del motor y parte del editor. En primer lugar, realicé una investigación sobre las posibles formas de adaptar el motor para que pudiera ser multiplataforma, realizando unas primeras pruebas técnicas y de concepto, y acabé decantándome por utilizar \texttt{SDL} como núcleo central del motor, principalmente alentado por todo el soporte multiplataforma que este aporta.

\medskip

Una vez concluido esto, y con una versión básica de \comillas{juego} funcional tanto en Android como en \textit{desktop}, desarrollamos un pequeño programa que mostraba un rectángulo de colores en pantalla, cuyos atributos (tamaño y color) se leían de un archivo \texttt{.lua}. 

A continuación, empecé junto a Miguel Curros García la fase de diseño del motor. En este paso intentamos dejar bien definidas las diferentes partes que íbamos a necesitar, así como la forma en que se comunicarían entre sí. Para ello, generamos un diagrama UML de los diferentes componentes que íbamos a utilizar y lo separamos en diferentes niveles de abstracción, dejando clara la separación entre los componentes básicos del motor y los componentes de \textit{gameplay} que íbamos a necesitar para poder implementar todas las funcionalidades específicas que queríamos para los juegos que ofrece nuestro editor. 

Una vez cerrado este diseño, comenzamos con la implementación. Yo me centré en toda la parte de \texttt{Render}, \texttt{Core} (sistema de Entidades y Componentes) e \texttt{Input}. En primer lugar, desarrollé el \textit{core} del motor: la estructura de Entidades y Componentes, todas ellas organizadas en escenas y gestionadas por un \texttt{SceneManager}. Creé cada una de las partes del ciclo de vida de estos elementos y monté toda la estructura para que pudieran ser cargados a partir de datos leídos de un archivo \texttt{Lua}. Para esto, nos apoyamos en \texttt{ComponentFactory} y en la clase \texttt{ComponentTemplate}, que, a través de una macro y una estructura de plantilla, aceleraba mucho el proceso de declarar nuevos componentes. A continuación, creé la clase \texttt{ComponentData}, que utilizaríamos para poder parametrizar los componentes que declarásemos en Lua. Finalmente, creé unos métodos en el \texttt{SceneManager} para poder instanciar entidades y escenas a partir de unos \textit{blueprints} creados leyendo los archivos Lua correspondientes. 

Una vez terminado esto, pasamos a la parte del \textit{renderizado}. Aquí, generalicé los métodos de pintado de \texttt{SDL} y creé un bucle de \textit{clear}, \textit{present} y \textit{render}, que se llama desde el bucle principal. Para acceder a las funciones de pintado, creé una clase virtual \texttt{RenderComponent} que implementa el método \texttt{render(RenderManager*)}. A continuación, creé todos los componentes básicos de \textit{renderizado} que íbamos a necesitar (\texttt{Camera}, \texttt{Rectangle}, \texttt{Text}, \texttt{SpriteRenderer}, \texttt{Animator}), así como los recursos que estos iban a utilizar (\texttt{Sprite}, \texttt{Animation}, \texttt{Font}, \texttt{Color}). Todo esto quedó integrado con el ciclo de vida que creé en el apartado de \texttt{Core}, para permitir su inicialización parametrizada desde Lua. 

Finalmente, hice una implementación sencilla de un sistema de \textit{input} al que se pudiera acceder desde los componentes. Dado que, por diseño, solo íbamos a utilizar clics/\textit{touchs} para mantener de una forma más simple el soporte multiplataforma, generalicé estos en un \texttt{struct}. Adicionalmente, creé el componente \texttt{Button}, al que se le podía asignar una función de Lua que se llamaría con unos parámetros preestablecidos al detectar un \textit{input} en su área. 

Con esto y las aportaciones de Miguel, dimos por terminado el motor genérico y pasamos a implementar los elementos específicos de \textit{gameplay}. Aquí desarrollé un sistema de diálogo formado principalmente por dos componentes: un gestor de \textit{textboxes}, que mostraba texto poco a poco y esperaba el \textit{input} del usuario, y unas opciones compuestas por varios botones con posibles respuestas por parte del usuario. Lo siguiente fue crear un sistema de movimiento basado en A*, que utilizarían tanto los NPC como el jugador a partir de un input de tipo \textit{point and click}; también añadí la opción de aplicar animaciones como parámetro a este movimiento. Una vez cerrado esto, creé un gestor de mundo cuya función sería controlar qué mapas están activos en escena en cada momento, instanciando nuevos en caso necesario a partir de los \textit{blueprints}. Todos estos sistemas serían la base del \textit{gameplay} de \textit{overworld} que ofrece nuestro editor, combinado con el sistema de eventos. 

Con esto listo, pasamos al desarrollo del editor, que ya tenía una base implementada por Alejandro Massó. Aquí me centré, en primer lugar, en generar un sistema de traducción que convirtiera los datos generados por el editor en datos que siguieran la estructura del motor (entidades, componentes, escenas, \textit{sprites} y animaciones), así como en la generación de otros archivos de configuración. A partir de esto, concreté un proceso de \textit{build} que se encargaría de obtener los binarios precompilados del motor en formato ejecutable y los combinaría con los \textit{assets} del usuario y los archivos de datos en formato motor, previamente traducidos, para obtener el producto final: el juego, contenido en un único directorio denominado \texttt{Build}.  

Posteriormente, implementé otras funcionalidades del editor, como el inspector de objetos, la pestaña de conexiones entre mapas, la pestaña de configuración del jugador y la de configuración general. Para esta última, quise implementar una previsualización para el texto del jugador, por lo que tuve que gestionar una carga asíncrona de fuentes con \texttt{DearImGui}. Todas estas funcionalidades supusieron también una ampliación del sistema de traducción y \textit{build} previamente comentado.  

Por último, me dediqué a realizar pruebas con usuarios y corregir los problemas que encontrábamos en estas. Finalmente, adapté la funcionalidad de \textit{build} para que también fuese posible generar una APK lista para usar en dispositivos Android.

\section*{Alejandro Massó Martínez}
Mi contribución al Proyecto se ha basado en la investigación, diseño y desarrollo del editor, la creación del \textit{toolchain} de generación de dependencias y compilación de los distintos subproyectos, y la redacción de esta memoria.

\medskip

%En primer lugar, hice una exhaustiva investigación acerca de 


%
% Bibliografía
%
% Si el TFM se escribe en inglés, editar TeXiS/TeXiS_bib para cambiar el
% estilo de las referencias
%---------------------------------------------------------------------
%
%                      configBibliografia.tex
%
%---------------------------------------------------------------------
%
% bibliografia.tex
% Copyright 2009 Marco Antonio Gomez-Martin, Pedro Pablo Gomez-Martin
%
% This file belongs to the TeXiS manual, a LaTeX template for writting
% Thesis and other documents. The complete last TeXiS package can
% be obtained from http://gaia.fdi.ucm.es/projects/texis/
%
% Although the TeXiS template itself is distributed under the 
% conditions of the LaTeX Project Public License
% (http://www.latex-project.org/lppl.txt), the manual content
% uses the CC-BY-SA license that stays that you are free:
%
%    - to share & to copy, distribute and transmit the work
%    - to remix and to adapt the work
%
% under the following conditions:
%
%    - Attribution: you must attribute the work in the manner
%      specified by the author or licensor (but not in any way that
%      suggests that they endorse you or your use of the work).
%    - Share Alike: if you alter, transform, or build upon this
%      work, you may distribute the resulting work only under the
%      same, similar or a compatible license.
%
% The complete license is available in
% http://creativecommons.org/licenses/by-sa/3.0/legalcode
%
%---------------------------------------------------------------------
%
% Fichero  que  configura  los  parámetros  de  la  generación  de  la
% bibliografía.  Existen dos  parámetros configurables:  los ficheros
% .bib que se utilizan y la frase célebre que aparece justo antes de la
% primera referencia.
%
%---------------------------------------------------------------------


%%%%%%%%%%%%%%%%%%%%%%%%%%%%%%%%%%%%%%%%%%%%%%%%%%%%%%%%%%%%%%%%%%%%%%
% Definición de los ficheros .bib utilizados:
% \setBibFiles{<lista ficheros sin extension, separados por comas>}
% Nota:
% Es IMPORTANTE que los ficheros estén en la misma línea que
% el comando \setBibFiles. Si se desea utilizar varias líneas,
% terminarlas con una apertura de comentario.
%%%%%%%%%%%%%%%%%%%%%%%%%%%%%%%%%%%%%%%%%%%%%%%%%%%%%%%%%%%%%%%%%%%%%%
\setBibFiles{%
biblio%
}

%%%%%%%%%%%%%%%%%%%%%%%%%%%%%%%%%%%%%%%%%%%%%%%%%%%%%%%%%%%%%%%%%%%%%%
% Definición de la frase célebre para el capítulo de la
% bibliografía. Dentro normalmente se querrá hacer uso del entorno
% \begin{FraseCelebre}, que contendrá a su vez otros dos entornos,
% un \begin{Frase} y un \begin{Fuente}.
%
% Nota:
% Si no se quiere cita, se puede eliminar su definición (en la
% macro setCitaBibliografia{} ).
%%%%%%%%%%%%%%%%%%%%%%%%%%%%%%%%%%%%%%%%%%%%%%%%%%%%%%%%%%%%%%%%%%%%%%
\setCitaBibliografia{
\begin{FraseCelebre}
\begin{Frase}
  Y así, del mucho leer y del poco dormir, se le secó el celebro de
  manera que vino a perder el juicio.\\ 
  \textcolor{red}{(modificar en Cascaras$\backslash$bibliografia.tex)}
\end{Frase}
\begin{Fuente}
  Miguel de Cervantes Saavedra
\end{Fuente}
\end{FraseCelebre}
}

%%
%% Creamos la bibliografia
%%
\makeBib

% Variable local para emacs, para  que encuentre el fichero maestro de
% compilación y funcionen mejor algunas teclas rápidas de AucTeX

%%%
%%% Local Variables:
%%% mode: latex
%%% TeX-master: "../Tesis.tex"
%%% End:



% Apéndices
\appendix
\chapter{Título del Apéndice A}
\label{Appendix:Key1}

Los apéndices son secciones al final del documento en las que se agrega texto con el objetivo de ampliar los contenidos del documento principal.
\chapter{Título del Apéndice B}
\label{Appendix:Key2}

Se pueden añadir los apéndices que se consideren oportunos.
%\include{Apendices/appendixC}
%\include{...}
%\include{...}
%\include{...}
\backmatter



%
% Índice de palabras
%

% Sólo  la   generamos  si  está   declarada  \generaindice.  Consulta
% TeXiS.sty para más información.

% En realidad, el soporte para la generación de índices de palabras
% en TeXiS no está documentada en el manual, porque no ha sido usada
% "en producción". Por tanto, el fichero que genera el índice
% *no* se incluye aquí (está comentado). Consulta la documentación
% en TeXiS_pream.tex para más información.
\ifx\generaindice\undefined
\else
%%---------------------------------------------------------------------
%
%                        TeXiS_indice.tex
%
%---------------------------------------------------------------------
%
% TeXiS_indice.tex
% Copyright 2009 Marco Antonio Gomez-Martin, Pedro Pablo Gomez-Martin
%
% This file belongs to TeXiS, a LaTeX template for writting
% Thesis and other documents. The complete last TeXiS package can
% be obtained from http://gaia.fdi.ucm.es/projects/texis/
%
% This work may be distributed and/or modified under the
% conditions of the LaTeX Project Public License, either version 1.3
% of this license or (at your option) any later version.
% The latest version of this license is in
%   http://www.latex-project.org/lppl.txt
% and version 1.3 or later is part of all distributions of LaTeX
% version 2005/12/01 or later.
%
% This work has the LPPL maintenance status `maintained'.
% 
% The Current Maintainers of this work are Marco Antonio Gomez-Martin
% and Pedro Pablo Gomez-Martin
%
%---------------------------------------------------------------------
%
% Contiene  los  comandos  para  generar  el índice  de  palabras  del
% documento.
%
%---------------------------------------------------------------------
%
% NOTA IMPORTANTE: el  soporte en TeXiS para el  índice de palabras es
% embrionario, y  de hecho  ni siquiera se  describe en el  manual. Se
% proporciona  una infraestructura  básica (sin  terminar)  para ello,
% pero  no ha  sido usada  "en producción".  De hecho,  a pesar  de la
% existencia de  este fichero, *no* se incluye  en Tesis.tex. Consulta
% la documentación en TeXiS_pream.tex para más información.
%
%---------------------------------------------------------------------


% Si se  va a generar  la tabla de  contenidos (el índice  habitual) y
% también vamos a  generar el índice de palabras  (ambas decisiones se
% toman en  función de  la definición  o no de  un par  de constantes,
% puedes consultar modo.tex para más información), entonces metemos en
% la tabla de contenidos una  entrada para marcar la página donde está
% el índice de palabras.

\ifx\generatoc\undefined
\else
   \addcontentsline{toc}{chapter}{\indexname}
\fi


% Generamos el índice
\printindex

% Variable local para emacs, para  que encuentre el fichero maestro de
% compilación y funcionen mejor algunas teclas rápidas de AucTeX

%%%
%%% Local Variables:
%%% mode: latex
%%% TeX-master: "./tesis.tex"
%%% End:

\fi

%
% Lista de acrónimos
%

% Sólo  lo  generamos  si  está declarada  \generaacronimos.  Consulta
% TeXiS.sty para más información.


\ifx\generaacronimos\undefined
\else
%---------------------------------------------------------------------
%
%                        TeXiS_acron.tex
%
%---------------------------------------------------------------------
%
% TeXiS_acron.tex
% Copyright 2009 Marco Antonio Gomez-Martin, Pedro Pablo Gomez-Martin
%
% This file belongs to TeXiS, a LaTeX template for writting
% Thesis and other documents. The complete last TeXiS package can
% be obtained from http://gaia.fdi.ucm.es/projects/texis/
%
% This work may be distributed and/or modified under the
% conditions of the LaTeX Project Public License, either version 1.3
% of this license or (at your option) any later version.
% The latest version of this license is in
%   http://www.latex-project.org/lppl.txt
% and version 1.3 or later is part of all distributions of LaTeX
% version 2005/12/01 or later.
%
% This work has the LPPL maintenance status `maintained'.
% 
% The Current Maintainers of this work are Marco Antonio Gomez-Martin
% and Pedro Pablo Gomez-Martin
%
%---------------------------------------------------------------------
%
% Contiene  los  comandos  para  generar  el listado de acrónimos
% documento.
%
%---------------------------------------------------------------------
%
% NOTA IMPORTANTE:  para que la  generación de acrónimos  funcione, al
% menos  debe  existir  un  acrónimo   en  el  documento.  Si  no,  la
% compilación  del   fichero  LaTeX  falla  con   un  error  "extraño"
% (indicando  que  quizá  falte  un \item).   Consulta  el  comentario
% referente al paquete glosstex en TeXiS_pream.tex.
%
%---------------------------------------------------------------------


% Redefinimos a español  el título de la lista  de acrónimos (Babel no
% lo hace por nosotros esta vez)

\def\listacronymname{Lista de acrónimos}

% Para el glosario:
% \def\glosarryname{Glosario}

% Si se  va a generar  la tabla de  contenidos (el índice  habitual) y
% también vamos a  generar la lista de acrónimos  (ambas decisiones se
% toman en  función de  la definición  o no de  un par  de constantes,
% puedes consultar config.tex  para más información), entonces metemos
% en la  tabla de contenidos una  entrada para marcar  la página donde
% está el índice de palabras.

\ifx\generatoc\undefined
\else
   \addcontentsline{toc}{chapter}{\listacronymname}
\fi


% Generamos la lista de acrónimos (en realidad el índice asociado a la
% lista "acr" de GlossTeX)

\printglosstex(acr)

% Variable local para emacs, para  que encuentre el fichero maestro de
% compilación y funcionen mejor algunas teclas rápidas de AucTeX

%%%
%%% Local Variables:
%%% mode: latex
%%% TeX-master: "../Tesis.tex"
%%% End:

\fi

%
% Final
%
%---------------------------------------------------------------------
%
%                      fin.tex
%
%---------------------------------------------------------------------
%
% fin.tex
% Copyright 2009 Marco Antonio Gomez-Martin, Pedro Pablo Gomez-Martin
%
% This file belongs to the TeXiS manual, a LaTeX template for writting
% Thesis and other documents. The complete last TeXiS package can
% be obtained from http://gaia.fdi.ucm.es/projects/texis/
%
% Although the TeXiS template itself is distributed under the 
% conditions of the LaTeX Project Public License
% (http://www.latex-project.org/lppl.txt), the manual content
% uses the CC-BY-SA license that stays that you are free:
%
%    - to share & to copy, distribute and transmit the work
%    - to remix and to adapt the work
%
% under the following conditions:
%
%    - Attribution: you must attribute the work in the manner
%      specified by the author or licensor (but not in any way that
%      suggests that they endorse you or your use of the work).
%    - Share Alike: if you alter, transform, or build upon this
%      work, you may distribute the resulting work only under the
%      same, similar or a compatible license.
%
% The complete license is available in
% http://creativecommons.org/licenses/by-sa/3.0/legalcode
%
%---------------------------------------------------------------------
%
% Contiene la última página
%
%---------------------------------------------------------------------


% Ponemos el marcador en el PDF
\ifpdf
   \pdfbookmark{Fin}{fin}
\fi

\thispagestyle{empty}\mbox{}

Este texto se puede encontrar en el fichero Cascaras/fin.tex. Si deseas eliminarlo, basta con comentar la línea correspondiente al final del fichero TFGTeXiS.tex.

\vspace*{4cm}

\small

\hfill \emph{--¿Qué te parece desto, Sancho? -- Dijo Don Quijote --}

\hfill \emph{Bien podrán los encantadores quitarme la ventura,}

\hfill \emph{pero el esfuerzo y el ánimo, será imposible.}

\hfill 

\hfill \emph{Segunda parte del Ingenioso Caballero} 

\hfill \emph{Don Quijote de la Mancha}

\hfill \emph{Miguel de Cervantes}

\vfill%space*{4cm}

\hfill \emph{--Buena está -- dijo Sancho --; fírmela vuestra merced.}

\hfill \emph{--No es menester firmarla -- dijo Don Quijote--,}

\hfill \emph{sino solamente poner mi rúbrica.}

\hfill 

\hfill \emph{Primera parte del Ingenioso Caballero} 

\hfill \emph{Don Quijote de la Mancha}

\hfill \emph{Miguel de Cervantes}


\newpage
\thispagestyle{empty}\mbox{}

\newpage

% Variable local para emacs, para  que encuentre el fichero maestro de
% compilación y funcionen mejor algunas teclas rápidas de AucTeX

%%%
%%% Local Variables:
%%% mode: latex
%%% TeX-master: "../Tesis.tex"
%%% End:

%\end{otherlanguage}
\end{document}
